\documentclass[12pt]{article}
\usepackage[utf8]{inputenc}
\usepackage[T1]{fontenc}
\usepackage[margin=1in]{geometry}
\usepackage{setspace}
\usepackage{xcolor}
\usepackage[colorlinks=true, urlcolor=blue!60!black, linkcolor=blue!60!black]{hyperref}
\setlength{\parindent}{0em}
\setlength{\parskip}{0.45em}
\renewcommand{\baselinestretch}{1.0}
\begin{document}
\begin{center}
{\Large \textit{Pierce v. Society of Sisters}, 268 U.S. 510 (1925)\par}
\vspace{0.3em}
{\normalsize Case Brief --- One-Page Edition\par}
\end{center}

\textbf{Memory Jogger.} A state may not ``standardize its children'' by abolishing private schools; parents have a fundamental liberty to direct their children's education (\href{https://supreme.justia.com/cases/federal/us/268/510/}{Justia}).

\textbf{Facts.} Amid post-war nativism and anti-Catholic sentiment backed by the Ku Klux Klan, Oregon voters in 1922 required all children aged eight to sixteen to attend public school, effectively outlawing private and parochial primary education. The Society of Sisters, a Catholic teaching order, and Hill Military Academy sued as parents began withdrawing students before the law took effect.

\textbf{Procedural History.} A three-judge district court enjoined enforcement; Oregon appealed directly.

\textbf{Votes.} Affirmed, 9--0, Justice McReynolds writing; no separate opinions.

\textbf{Holding.} Under \textit{Meyer v. Nebraska}, compelling public-school attendance unreasonably interferes with parents' liberty to direct their children's upbringing; private schools may sue before the effective date to protect threatened business and property.

\textbf{Reasoning.} The state may require attendance at some school and reasonably regulate all schools (technically dicta), but it may not destroy useful private schools without any emergency or reasonable relation to a legitimate purpose. Corporations lack Fourteenth Amendment liberty, but the schools' business and property were threatened by unconstitutional compulsion of their patrons, and present losses made the suits ripe.

\textbf{Examples.} \emph{Same-Side:} a ban on private middle schools. \emph{Opposite-Side:} fire codes and core-subject requirements for private schools. \emph{Fence-Sitter:} a secular-equivalency mandate threatening yeshivas.

\textbf{Critique.} Born in the \textit{Lochner} era, it rests on unenumerated liberty, and child-rights scholars say it treats children as parental property.

\textbf{Key Quotation.} ``Under the doctrine of \textit{Meyer v. Nebraska}, 262 U.\,S. 390, we think it entirely plain that the Act of 1922 unreasonably interferes with the liberty of parents and guardians to direct the upbringing and education of children under their control.'' --- Justice McReynolds, 268 U.S. at 534.

\end{document}
